\section{Introduction}
Structural variants are a major source of oncogenic activation, tumor
suppressor loss, enhancer hijacking, and genome reorganization. In
pediatric tumors, however, broad pan-cancer summaries of structural
variation have historically lagged behind point-mutation catalogs, which
made it harder to compare disease classes within a unified framework.

Earlier pediatric pan-cancer genomics efforts and whole-cancer
structural-variation surveys provided important context for that gap by
outlining recurrent genomic alterations across childhood cancers and by
formalizing major rearrangement patterns across human malignancies
\supercite{Grobner_2018,Ma_2018,Li_2020}.

The public study by Greenhalgh and colleagues addressed that gap by
surveying pediatric and adult whole-genome cohorts side by side and by
focusing on the biological consequences of rearrangements instead of
only their raw counts \supercite{Greenhalgh_2026}. Their work is especially useful as
a formatting example because it combines large-cohort genomics, mutation
signatures, and clinically relevant language in a conventional IMRaD
structure.

This project does not reproduce the journal article verbatim. Instead,
it uses publicly accessible metadata and abstract-level findings to
create a clean demonstration manuscript for the \texttt{bensz-paper}
package. The goal is to give repository users a ready-to-render SCI
template with both PDF and Word outputs while keeping the sample text
openly shareable.
